% ---
% id: protocol-compiler/rem:completeness
% kind: rem
% title: "The two directions of the validator"
% summary: >-
%   Soundness keeps an ill-formed declaration from being emitted; completeness,
%   which holds for every target condition but feasibility, makes a refusal
%   name a real defect. A refusal at V7 means only that no verified feasible
%   state was shown.
% prominence: standard
% status: accepted
% uses: [thm:validator-soundness, thm:feasibility]
% mentions: [req:witness-refusal-meaning, req:feasibility-witness]
% assumes: []
% refs: []
% provenance: newly-derived
% lean:
% ---

\begin{remark}
\label{rem:completeness}
The two directions do different jobs and it matters which one is available
where. \emph{Soundness} --- Theorem~\ref{thm:validator-soundness}\,(i), and it holds for all
eight conditions --- is the safety direction: without it the compiler could
emit a declaration that passes every check and is nonetheless ill-formed, which
is the failure this module exists to prevent, since a downstream engine
consumes D(i)--D(viii) as hypotheses without re-deriving them. \emph{Completeness}
is what makes a \emph{refusal} informative: where it holds, a well-formed
declaration is never rejected, so a refusal names a real defect and not a
limitation of the checker. It holds for D(i)--D(vi) and D(viii)
(Theorem~\ref{thm:validator-soundness}\,(ii)); both directions of each of those pairs are
proved, not just the one that keeps bad declarations out. Both directions
are read on the declared numbers: a pass buys D(i)--D(viii) of the declared
numbers, which is what the consuming engine requires of them, and a
declaration well-formed of those numbers passes every check but V7.

Completeness does \emph{not} hold for D(vii), and that is a limitation of
this checker (Theorem~\ref{thm:validator-soundness}\,(iii)). \emph{Its scope is the third
branch of V7 and no wider}: a declaration with an empty constraint list
passes on branch (a), and one whose feasible set contains the default
candidate \eqref{eq:witness-empty} passes on branch (b), with nothing supplied
in either case. What is refused is a well-formed declaration with a non-empty
constraint list whose feasible set does not contain the default candidate and
for which no witness was supplied. Three consequences follow. A refusal at V7
means ``I was not shown a feasible state, and the one I tried is not one'',
never ``there is none'' (Requirement~\ref{req:witness-refusal-meaning}).
Whoever presents a declaration refused for that reason has a remedy inside
this interface --- supply $\omega^{\star}$
(Requirement~\ref{req:feasibility-witness}) --- and no other check has that
character. And the asymmetry is the deliberate one: the direction that is
missing costs a rejected good declaration, while the direction that is present
is the one whose absence would cost a consumed bad one.
\end{remark}
